\documentclass[12pt,a4paper]{exam}

% --------------------
% Packages
% --------------------
\usepackage[utf8]{inputenc}
\usepackage[T1]{fontenc}
\usepackage[french]{babel}
\usepackage{amsmath,amssymb}
\usepackage{geometry}
\usepackage{pgfplots}
\usepackage{tkz-tab}
\pgfplotsset{compat=1.18}
\geometry{margin=2cm}

% --------------------
% En-tête / pied
% --------------------
\header{}{Corrigé – DS de mathématiques}{1ère STMG}
\footer{}{\thepage}{}

% --------------------
\begin{document}

\begin{center}
    \Large\textbf{Corrigé du devoir surveillé\\Fonctions polynômes de degré 2}
\end{center}

\vspace{0.5cm}

% ======================================================
\noindent
\textbf{Exercice 1 : Étude d’une fonction \ (/5 points)}

\vspace{5pt}

\begin{questions}

\question \textbf{Valeurs de $f(0)$ et $f(1)$}

\begin{itemize}
    \item $f(0)$ correspond à l’ordonnée à l’origine : on lit sur le graphique  
    \[
    f(0) = -3
    \]
    \item Pour $x=1$, on lit graphiquement :
    \[
    f(1) = 0
    \]
\end{itemize}

\question \textbf{Détermination de $a$ et $c$}

On sait que :
\[
f(x) = ax^2 + c
\]

\begin{itemize}
    \item Comme $f(0) = c$ (la valeur sur les ordonnées où passe la 
    courbe), on obtient :
    \[
    c = -3
    \]
    \item Comme $f(1) = a \times 1^2 + c = a - 3$ et que $f(1) = 0$ (graphiquement), on a :
    \[
    a - 3 = 0 \quad \Rightarrow \quad a = 3
    \]
\end{itemize}

\[
\boxed{f(x)=3x^2-3}
\]

\end{questions}

\vspace{0.5cm}

% ======================================================
\noindent
\textbf{Exercice 2 : Étudier et tracer \ (/8 points)}

\vspace{0.3cm}
On considère la fonction :
\[
f(x) = -2(x+1)(x-5)
\]

\begin{questions}

\question \textbf{Caractéristiques du polynôme}

\begin{parts}

\part \textbf{Racines}

\[
x_1=-1 \quad \text{et} \quad x_2=5
\]

\part \textbf{Axe de symétrie}

L’axe est la moyenne des racines :
\[
x=\dfrac{x_1 + x_2}{2} = \frac{-1+5}{2}=2
\]

\part \textbf{Sommet}

Coordonnées du sommet : $\left(\dfrac{x_1+x_2}{2}\  ; \ f(\dfrac{x_1+x_2}{2})\right) = (2, f(2))$

Avec :

\[
f(2)=-2(3)(-3)=18
\]
Donc :
\[
S(2;18)
\]

\part \textbf{Nature du sommet}

Le coefficient directeur $a$ est négatif ($a = -2<0$), donc la parabole est tournée vers le bas.
\[
\boxed{\text{Le sommet est un maximum}}
\]

\end{parts}

\question \textbf{Allure de la courbe}

La courbe est une parabole :
\begin{itemize}
    \item tournée vers le bas,
    \item coupant l’axe des abscisses en $-1$ et $5$,
    \item de sommet $S(2;18)$.
\end{itemize}

\vspace{10pt}

\question \textbf{Tableau de variations}

\[
\begin{tikzpicture}
   \tkzTabInit{$x$ / 1 , $f(x)$ / 1.5}{$-\infty$, $2$, $+\infty$}
   \tkzTabVar{-/ , +/ 18, -/}
\end{tikzpicture}
\]

\end{questions}

\vspace{0.5cm}

% ======================================================
\noindent
\textbf{Exercice 3 : Problème \ (/7 points)}

\vspace{0.3cm}
On considère la fonction :
\[
f(x) = -2x^2 + 2x + 4
\]

\begin{questions}

\question \textbf{Montrer que 2 est une racine}

\[
f(2) = -2 \times 2^2 + 2 \times 2 + 4 = -2 \times 4 + 4 + 4 = -8 + 8 = 0
\]
Donc $2$ est bien une racine.

\question \textbf{Calcul de $f(-1)$}

\[
f(-1) = -2 \times (-1)^2 + 2 \times (-1) + 4 = - 2 - 2 + 4 = 0
\]
Donc $-1$ est également une racine.

\question \textbf{Forme factorisée}

Les racines sont $-1$ et $2$ :
\[
\boxed{f(x) = -2(x+1)(x-2)}
\]

\question \textbf{Tableau de signes}

\[
\begin{tikzpicture}
   \tkzTabInit{$x$ / 1 , $f(x)$ / 1}{$-\infty$, $-1$, $2$, $+\infty$}
   \tkzTabLine{, -, z, +, z, -, }
\end{tikzpicture}
\]

\question \textbf{Résolution de l’inéquation}

\[
f(x) \leq 0 \iff x \leq -1 \quad \text{ou} \quad x \geq 2
\]

\[
\boxed{S = \ \left]-\infty;-1\right] \cup \left[2;+\infty\right[}
\]

\end{questions}

\vspace{0.5cm}

\begin{center}
    \textit{Fin du corrigé}
\end{center}

\end{document}
